% GENERATED FILE. Edit canonical lessons, then run npm run generate:books.
% canonical-source-hash: fnv1a64:cce51d46d4687f8f
% canonical-lessons: JA-C533-yarigai, JA-C533-ikigai, JA-C533-netchuu-suru, JA-C533-uchikomu, JA-C533-joutatsu-suru, JA-R533-review-528-533, JA-R533-review-final-533

\chapter{Worth the Effort}
\label{ch:ja-worth-the-effort}
\hlchaptermodality{533}

\begin{chapteropening}
\textbf{By the end of this chapter:} \emph{I can say what would make an effort worthwhile.}

Complete the last lesson of chapter 533: I can say what would make an effort worthwhile.
\end{chapteropening}

\section[\texorpdfstring{\ja{やりがい} (yarigai)}{やりがい (yarigai)}]{\ja{やりがい} (yarigai) — a sense of reward; satisfaction (in a task)}

\label{lesson:JA-C533-yarigai}



\begin{quote}
Before the new one: say the Japanese for advantageous, then the Japanese for surprising.
\end{quote}

\subsection*{You'll want to know: \ja{やりがい}}
\textbf{\ja{やりがい}} — \emph{yarigai} — \textquotedblleft{}a sense of reward; satisfaction (in a task)\textquotedblright{}.

\begin{quote}
\textbf{\ja{やりがいは} \ja{しごとに} \ja{よります。}} — \emph{yarigai wa shigoto ni yorimasu} — \textquotedblleft{}How rewarding it is depends on the job.\textquotedblright{}
\end{quote}

\begin{cousinweb}
\textbf{\ja{やりがい}} (\emph{yarigai}) is a native Japanese word (wago), inherited rather than borrowed from Chinese, and usually written in kana; its older origin is not settled (JMdict).
\end{cousinweb}

\begin{grammarlens}[title={What you've built}]
One more word for saying what happens if.
\end{grammarlens}

\subsection*{Your turn}
\begin{itemize}
\raggedright
  \item \emph{Say it:} \emph{yarigai}
  \item \emph{Say it:} \emph{yarigai}, once more
  \item \emph{Say it:} \emph{yarigai wa shigoto ni yorimasu} — How rewarding it is depends on the job.
  \item \emph{Recall:} say the Japanese for advantageous, then the Japanese for surprising, then say \emph{yarigai} again
\end{itemize}

\begin{tcolorbox}[breakable,colback=teal!4,colframe=teal!35!black,title={Before you move on}]
What does \ja{やりがい} mean? (\textquotedblleft{}A sense of reward; satisfaction (in a task)\textquotedblright{}.) Say it once more.
\end{tcolorbox}
\section[\texorpdfstring{\ja{いきがい} (ikigai)}{いきがい (ikigai)}]{\ja{いきがい} (ikigai) — a purpose in life}

\label{lesson:JA-C533-ikigai}



\begin{quote}
Before the new one: say the Japanese for surprising, then the Japanese for a sense of reward.
\end{quote}

\subsection*{You'll want to know: \ja{いきがい}}
\textbf{\ja{いきがい}} — \emph{ikigai} — \textquotedblleft{}a purpose in life\textquotedblright{}.

\begin{quote}
\textbf{\ja{もし} \ja{いきがいが} \ja{みつかったら、まいにち} \ja{たのしいでしょう。}} — \emph{moshi ikigai ga mitsukattara, mainichi tanoshii deshō} — \textquotedblleft{}If you find a purpose in life, every day will surely be enjoyable.\textquotedblright{}
\end{quote}

\begin{cousinweb}
\textbf{\ja{いきがい}} (\emph{ikigai}) is a native Japanese word (wago), inherited rather than borrowed from Chinese; its kanji was attached later as a kun'yomi, and its older origin is not settled (JMdict, KANJIDIC2).
\end{cousinweb}

\begin{grammarlens}[title={What you've built}]
One more word for saying what happens if.
\end{grammarlens}

\subsection*{Your turn}
\begin{itemize}
\raggedright
  \item \emph{Say it:} \emph{ikigai}
  \item \emph{Say it:} \emph{ikigai}, once more
  \item \emph{Say it:} \emph{moshi ikigai ga mitsukattara, mainichi tanoshii deshō} — If you find a purpose in life, every day will surely be enjoyable.
  \item \emph{Recall:} say the Japanese for surprising, then the Japanese for a sense of reward, then say \emph{ikigai} again
\end{itemize}

\begin{tcolorbox}[breakable,colback=teal!4,colframe=teal!35!black,title={Before you move on}]
What does \ja{いきがい} mean? (\textquotedblleft{}A purpose in life\textquotedblright{}.) Say it once more.
\end{tcolorbox}
\section[\texorpdfstring{\ja{ねっちゅうする} (netchū suru)}{ねっちゅうする (netchū suru)}]{\ja{ねっちゅうする} (netchū suru) — to be absorbed (in)}

\label{lesson:JA-C533-netchuu-suru}



\begin{quote}
Before the new one: say the Japanese for a sense of reward, then the Japanese for a purpose in life.
\end{quote}

\subsection*{You'll want to know: \ja{ねっちゅうする}}
\textbf{\ja{ねっちゅうする}} — \emph{netchū suru} — \textquotedblleft{}to be absorbed (in)\textquotedblright{}.

\begin{quote}
\textbf{\ja{もし} \ja{しゅみに} \ja{ねっちゅうしたら、あには} \ja{じかんを} \ja{わすれるでしょう。}} — \emph{moshi shumi ni netchū shitara, ani wa jikan o wasureru deshō} — \textquotedblleft{}If my elder brother gets absorbed in a hobby, he will forget the time.\textquotedblright{}
\end{quote}

Polite \textbf{\ja{ねっちゅうします}}, \emph{netchūshimasu}; a \ja{する} verb. Past (polite): \textbf{\ja{おとうとが} \ja{しょうぎに} \ja{ねっちゅうしました。}}, \emph{otōto ga shōgi ni netchūshimashita}.

\begin{cousinweb}
\textbf{\ja{ねっちゅうする}} (\emph{netchū suru}) is the Sino-Japanese (kango) noun \emph{netchū} with the native verb \emph{suru}, \textquotedblleft{}to do\textquotedblright{}: two Chinese characters read with their on'yomi, readings borrowed from Chinese. They are \emph{ne}, \textquotedblleft{}heat\textquotedblright{} (Middle Chinese \emph{nyet}) and \emph{chū}, \textquotedblleft{}in\textquotedblright{} (\emph{trjuwng}) (KANJIDIC2; Middle Chinese in Baxter's transcription).
\end{cousinweb}

\begin{grammarlens}[title={What you've built}]
One more word for saying what happens if.
\end{grammarlens}

\subsection*{Your turn}
\begin{itemize}
\raggedright
  \item \emph{Say it:} \emph{netchū suru}
  \item \emph{Say it:} \emph{netchū suru}, once more
  \item \emph{Say it:} \emph{moshi shumi ni netchū shitara, ani wa jikan o wasureru deshō} — If my elder brother gets absorbed in a hobby, he will forget the time.
  \item \emph{Recall:} say the Japanese for a sense of reward, then the Japanese for a purpose in life, then say \emph{netchū suru} again
\end{itemize}

\begin{tcolorbox}[breakable,colback=teal!4,colframe=teal!35!black,title={Before you move on}]
What does \ja{ねっちゅうする} mean? (\textquotedblleft{}To be absorbed (in)\textquotedblright{}.) Say it once more.
\end{tcolorbox}
\section[\texorpdfstring{\ja{うちこむ} (uchikomu)}{うちこむ (uchikomu)}]{\ja{うちこむ} (uchikomu) — to throw oneself (into)}

\label{lesson:JA-C533-uchikomu}



\begin{quote}
Before the new one: say the Japanese for a purpose in life, then the Japanese for to be absorbed (in).
\end{quote}

\subsection*{You'll want to know: \ja{うちこむ}}
\textbf{\ja{うちこむ}} — \emph{uchikomu} — \textquotedblleft{}to throw oneself (into)\textquotedblright{}.

\begin{quote}
\textbf{\ja{もし} \ja{べんきょうに} \ja{うちこんだら、いもうとは} \ja{しけんに} \ja{うかるでしょう。}} — \emph{moshi benkyō ni uchikondara, imōto wa shiken ni ukaru deshō} — \textquotedblleft{}If my younger sister throws herself into her studies, she will pass the exam.\textquotedblright{}
\end{quote}

Polite \textbf{\ja{うちこみます}}, \emph{uchikomimasu}; a godan verb. Past (polite): \textbf{\ja{あねが} \ja{れんしゅうに} \ja{うちこみました。}}, \emph{ane ga renshū ni uchikomimashita}.

\begin{cousinweb}
\textbf{\ja{うちこむ}} (\emph{uchikomu}) is a native Japanese word (wago), inherited rather than borrowed from Chinese; its kanji was attached later as a kun'yomi, and its older origin is not settled (JMdict, KANJIDIC2).
\end{cousinweb}

\begin{grammarlens}[title={What you've built}]
One more word for saying what happens if.
\end{grammarlens}

\subsection*{Your turn}
\begin{itemize}
\raggedright
  \item \emph{Say it:} \emph{uchikomu}
  \item \emph{Say it:} \emph{uchikomu}, once more
  \item \emph{Say it:} \emph{moshi benkyō ni uchikondara, imōto wa shiken ni ukaru deshō} — If my younger sister throws herself into her studies, she will pass the exam.
  \item \emph{Recall:} say the Japanese for a purpose in life, then the Japanese for to be absorbed (in), then say \emph{uchikomu} again
\end{itemize}

\begin{tcolorbox}[breakable,colback=teal!4,colframe=teal!35!black,title={Before you move on}]
What does \ja{うちこむ} mean? (\textquotedblleft{}To throw oneself (into)\textquotedblright{}.) Say it once more.
\end{tcolorbox}
\section[\texorpdfstring{\ja{じょうたつする} (jōtatsu suru)}{じょうたつする (jōtatsu suru)}]{\ja{じょうたつする} (jōtatsu suru) — to improve (in a skill)}

\label{lesson:JA-C533-joutatsu-suru}



\begin{quote}
Before the new one: say the Japanese for to be absorbed (in), then the Japanese for to throw oneself (into).
\end{quote}

\subsection*{You'll want to know: \ja{じょうたつする}}
\textbf{\ja{じょうたつする}} — \emph{jōtatsu suru} — \textquotedblleft{}to improve (in a skill)\textquotedblright{}.

\begin{quote}
\textbf{\ja{もし} \ja{まいにち} \ja{れんしゅうしたら、あねは} \ja{じょうたつするでしょう。}} — \emph{moshi mainichi renshū shitara, ane wa jōtatsu suru deshō} — \textquotedblleft{}If my elder sister practises every day, she will get better.\textquotedblright{}
\end{quote}

Polite \textbf{\ja{じょうたつします}}, \emph{jōtatsushimasu}; a \ja{する} verb. Past (polite): \textbf{\ja{あにの} \ja{日本語}\ja{が} \ja{じょうたつしました。}}, \emph{ani no nihongo ga jōtatsushimashita}.

\begin{cousinweb}
\textbf{\ja{じょうたつする}} (\emph{jōtatsu suru}) is the Sino-Japanese (kango) noun \emph{jōtatsu} with the native verb \emph{suru}, \textquotedblleft{}to do\textquotedblright{}: two Chinese characters read with their on'yomi, readings borrowed from Chinese. They are \emph{jō}, \textquotedblleft{}above\textquotedblright{} (Middle Chinese \emph{dzyangH}) and \emph{tatsu}, \textquotedblleft{}accomplished\textquotedblright{} (\emph{dat}) (KANJIDIC2; Middle Chinese in Baxter's transcription).
\end{cousinweb}

\begin{grammarlens}[title={What you've built}]
One more word for saying what happens if.
\end{grammarlens}

\subsection*{Your turn}
\begin{itemize}
\raggedright
  \item \emph{Say it:} \emph{jōtatsu suru}
  \item \emph{Say it:} \emph{jōtatsu suru}, once more
  \item \emph{Say it:} \emph{moshi mainichi renshū shitara, ane wa jōtatsu suru deshō} — If my elder sister practises every day, she will get better.
  \item \emph{Recall:} say the Japanese for to be absorbed (in), then the Japanese for to throw oneself (into), then say \emph{jōtatsu suru} again
\end{itemize}

\begin{tcolorbox}[breakable,colback=teal!4,colframe=teal!35!black,title={Before you move on}]
What does \ja{じょうたつする} mean? (\textquotedblleft{}To improve (in a skill)\textquotedblright{}.) Say it once more.
\end{tcolorbox}
\section[(dialogue)]{Looking back: from This or That to Worth the Effort}

\label{lesson:JA-R533-review-528-533}



\begin{quote}
Say the Japanese for or perhaps and or.
\end{quote}

\subsection*{Your turn}
Say these aloud:

\begin{itemize}
\raggedright
  \item or perhaps, or, or (in a question), or else, either one
  \item for some reason, by the way, anyway, would sooner, as one would expect
  \item fairly, at the very least, nearly, about, (not) at all
  \item on purpose, without meaning to, recklessly, desirable, somehow
  \item to limit, a limit, extra, advantageous, surprising
  \item a sense of reward, a purpose in life, to be absorbed (in), to throw oneself (into), to improve (in a skill)
\end{itemize}

\begin{tcolorbox}[breakable,colback=teal!4,colframe=teal!35!black,title={Before you move on}]
Or perhaps? (\textbf{\ja{あるいは}}, \emph{aruiwa}.) Or? (\textbf{\ja{または}}, \emph{matawa}.) Or (in a question)? (\textbf{\ja{それとも}}, \emph{soretomo}.) Or else? (\textbf{\ja{もしくは}}, \emph{moshikuwa}.) Either one? (\textbf{\ja{いずれか}}, \emph{izureka}.) For some reason? (\textbf{\ja{なぜか}}, \emph{nazeka}.) By the way? (\textbf{\ja{ところで}}, \emph{tokorode}.) Anyway? (\textbf{\ja{どうせ}}, \emph{dōse}.) Would sooner? (\textbf{\ja{いっそ}}, \emph{isso}.) As one would expect? (\textbf{\ja{さすが}}, \emph{sasuga}.) Fairly? (\textbf{\ja{かなり}}, \emph{kanari}.) At the very least? (\textbf{\ja{すくなくとも}}, \emph{sukunakutomo}.) Nearly? (\textbf{\ja{ほぼ}}, \emph{hobo}.) About? (\textbf{\ja{およそ}}, \emph{oyoso}.) (not) at all? (\textbf{\ja{まったく}}, \emph{mattaku}.) On purpose? (\textbf{\ja{わざと}}, \emph{wazato}.) Without meaning to? (\textbf{\ja{つい}}, \emph{tsui}.) Recklessly? (\textbf{\ja{むやみに}}, \emph{muyamini}.) Desirable? (\textbf{\ja{のぞましい}}, \emph{nozomashii}.) Somehow? (\textbf{\ja{なんとなく}}, \emph{nantonaku}.) To limit? (\textbf{\ja{かぎる}}, \emph{kagiru}.) A limit? (\textbf{\ja{せいげん}}, \emph{seigen}.) Extra? (\textbf{\ja{よぶん}}, \emph{yobun}.) Advantageous? (\textbf{\ja{ゆうり}}, \emph{yūri}.) Surprising? (\textbf{\ja{いがい}}, \emph{igai}.) A sense of reward? (\textbf{\ja{やりがい}}, \emph{yarigai}.) A purpose in life? (\textbf{\ja{いきがい}}, \emph{ikigai}.) To be absorbed (in)? (\textbf{\ja{ねっちゅうする}}, \emph{netchū suru}.) To throw oneself (into)? (\textbf{\ja{うちこむ}}, \emph{uchikomu}.) To improve (in a skill)? (\textbf{\ja{じょうたつする}}, \emph{jōtatsu suru}.)
\end{tcolorbox}
\section[(dialogue)]{Worth the Effort — once more}

\label{lesson:JA-R533-review-final-533}



\begin{quote}
Say the Japanese for a sense of reward and a purpose in life.
\end{quote}

\subsection*{Your turn}
\begin{itemize}
\raggedright
  \item \emph{Say it:} a sense of reward, a purpose in life, to be absorbed (in), to throw oneself (into), to improve (in a skill)
\end{itemize}

\begin{tcolorbox}[breakable,colback=teal!4,colframe=teal!35!black,title={Before you move on}]
A sense of reward? (\textbf{\ja{やりがい}}, \emph{yarigai}.) A purpose in life? (\textbf{\ja{いきがい}}, \emph{ikigai}.) To be absorbed (in)? (\textbf{\ja{ねっちゅうする}}, \emph{netchū suru}.) To throw oneself (into)? (\textbf{\ja{うちこむ}}, \emph{uchikomu}.) To improve (in a skill)? (\textbf{\ja{じょうたつする}}, \emph{jōtatsu suru}.)
\end{tcolorbox}
